\section{Experimental Design and Evaluation}
\label{sec:framework}

Section~\ref{sec:evidence-standard} defines the evidence required to convert controlled brain predictivity into a brain-specific training claim.
This section instantiates that evidence standard using distinct data, interventions, assays, comparators, endpoints, and inference units.
An estimand is the specific quantity or effect that an analysis is designed to estimate.
Table~\ref{tab:evidence-map} maps each conceptual stage to its thesis-facing scientific role and corresponding design tuple.
A stage may appear in more than one row when different datasets, response constructions, or population units define different estimands.
This section specifies the study design; Section~\ref{sec:results} reports the verdict for each role.

\begin{table}[H]
  \centering
  \footnotesize
  \setlength{\tabcolsep}{4pt}
  \renewcommand{\arraystretch}{1.08}
\caption{Scientific roles within the evidence chain.}
  \label{tab:evidence-map}
  \begin{tabularx}{\textwidth}{
    @{}
    >{\raggedright\arraybackslash}p{9mm}
    >{\raggedright\arraybackslash}p{32mm}
    >{\raggedright\arraybackslash}X
    >{\raggedright\arraybackslash}p{31mm}
    @{}
  }
    \toprule
    Stage & Scientific role & Primary comparison and endpoint & Inference unit \\
    \midrule
    1 & Controlled regional predictivity assay & Trained representations versus nuisance and architecture-matched untrained \ac{lm}; held-out ordinary and unique $R^2$ & Fixed sentence set; five contiguous folds \\
    \specialrule{0.15pt}{0.35ex}{0.35ex}
    1 & Controlled naturalistic voxelwise predictivity assay & Trained versus untrained representations beyond nuisance; held-out-story voxelwise unique $R^2$ & Story blocks within one deeply sampled participant \\
    \specialrule{0.15pt}{0.35ex}{0.35ex}
    2 & Quality-aware distillation-headroom analysis & Declared-layer brain predictivity before and after ordinary \ac{kd}, interpreted with held-out \ac{lm} quality & Model condition; training seed for stochastic arms \\
    \specialrule{0.15pt}{0.35ex}{0.35ex}
    3 & Exact synthetic-target measurability assay & Synthetic target versus nuisance and frozen row twin; held-out recorded-response unique $R^2$ & Participant on a fixed sentence set \\
    \specialrule{0.15pt}{0.35ex}{0.35ex}
    4 & Synthetic-target recovery assay & Target-trained versus \ac{kd} and target permutations; fresh held-out target recovery & Training seed \\
    \specialrule{0.15pt}{0.35ex}{0.35ex}
    5 & Retained-student movement audit & Target-trained versus \ac{kd}; retained parameter or representation change at comparable \ac{lm} quality & Training seed; fixed target sample for geometry \\
    \specialrule{0.15pt}{0.35ex}{0.35ex}
    6 & Recorded-response intervention, participant-averaged estimand & Correctly paired versus permuted response targets; fresh averaged-response unique $R^2$ & Shared stimulus fold; seeds are technical repetitions \\
    \specialrule{0.15pt}{0.35ex}{0.35ex}
    6 & Recorded-response intervention, participant-specific estimand & Correctly paired versus permuted response targets; fresh participant-response unique $R^2$ & Participant after fold and seed aggregation \\
    \specialrule{0.15pt}{0.35ex}{0.35ex}
    6 & Comparator-adequacy audit and attribution assessment & Synthetic brain-response versus text-derived auxiliary-control arm; matching checks and claimed outcome & Seed for trained models; fixed target sample for matching checks \\
    \specialrule{0.15pt}{0.35ex}{0.35ex}
    7 & Saved-student biological-transfer assay & Target-trained versus \ac{kd} and text-derived auxiliary-control arm; fresh participant-response unique $R^2$ & Participant \\
    \specialrule{0.15pt}{0.35ex}{0.35ex}
    8 & Bounded practical-utility evaluation & Brain-response-target arm versus route-appropriate quality-matched control; held-out language or task endpoint & Seed and task example, as declared \\
    \bottomrule
  \end{tabularx}
\end{table}

\subsection{Study Design, Datasets, and Data Partitioning}
\label{sec:data-targets}

The Tuckute condition-B data provide the fixed sentence-level dataset used for controlled brain predictivity, recorded-response training, and saved-student biological-transfer evaluation \autocite{tuckute2024}.
The dataset contains 1,000 isolated English sentences ordered by item identifier, recorded \ac{fmri} responses from 10 participants, and five fixed left-hemisphere language \acp{roi} (\texttt{AntTemp}, \texttt{IFG}, \texttt{IFGorb}, \texttt{MFG}, and \texttt{PostTemp}).
The participant-averaged assay uses \acp{uid} 848, 853, 865, 875, and 876.
Available responses are averaged cellwise, and inclusion requires a complete sentence-by-\ac{roi} matrix.
Participant-specific analyses exclude \ac{uid}~853 because of incomplete \ac{roi} coverage, leaving nine participants.
All Tuckute sentence analyses preserve item order and use five contiguous folds, each holding out 200 sentences while fitting on the remaining 800.
The released five-\ac{roi} reliability value, \(r_{\mathrm{NC}}=0.491\), provides descriptive correlation-scale context; raw variance-scale \(R^2\) remains the evaluation quantity \evd{E002}\evd{E004}\evd{E005}\evd{E008}\evd{E025}.

The LeBel corpus provides a naturalistic, voxelwise complement to the isolated-sentence, \ac{roi}-level Tuckute assay \autocite{lebel2023}.
The analysis uses participant UTS03, 20 stories divided into five held-out-story folds, and 11,442 voxels whose corrected split-half reliability exceeds 0.5 on a separate repeated story.
Holding out complete stories preserves within-story temporal dependence and tests prediction on unseen narratives, while the separate reliability story keeps evaluation responses outside voxel selection.
The resulting estimand concerns held-out-story prediction within one deeply sampled participant.
Stories and voxels provide nested measurements of that participant; population inference across people requires multiple independent participants.
Appendix~\ref{app:data-targets} specifies stimulus alignment, temporal resampling, response preprocessing, and delay construction, while Section~\ref{sec:measurement} describes the held-out readout protocol \evd{E006}.

Response construction determines the quantity being estimated.
Participant averaging can increase reliability by emphasizing the stimulus-locked response shared across selected participants, producing one group-level sentence-by-\ac{roi} matrix.
Scores against this matrix estimate predictivity of that shared response; population replication still depends on independent participants.
Participant-specific construction instead preserves one response matrix and one intervention effect for each valid participant, with folds and seeds aggregated within that person before population inference.
At the spatial level, each \ac{roi} coordinate summarizes activity across a predefined language region, whereas each voxel coordinate retains a finer-grained response measurement.
The Tuckute assays therefore estimate sentence-level predictivity over five regional summaries, while the LeBel assay estimates held-out-story predictivity over reliable voxels within one participant.
A response construction becomes a training target when it enters the auxiliary objective and an evaluation endpoint when a fresh held-out readout predicts it.
Section~\ref{sec:estimands} pairs each construction and role with its corresponding estimand and inference unit \evd{E002}\evd{E006}\evd{E008}\evd{E014}\evd{E017}\evd{E025}.

Each dataset has a fixed role in the evidence chain.
The recorded-response training family uses the five outer folds of the 1,000 Tuckute sentences: each arm trains on 800 sentences, and the remaining 200 provide the fresh recorded-response endpoint.
A separate held-out WikiText set measures \ac{lm} quality.
The synthetic-response family trains on 95,999 WikiText-103 sentences with fixed text-derived targets; its separate 1,999-sentence split supplies language-quality and held-out target-recovery endpoints.
The trained students are subsequently frozen before evaluation on the Tuckute responses used for exact-target measurability and biological transfer.
The LeBel stories remain confined to the independent naturalistic controlled-predictivity assay.
Fold-dependent transformations are fitted only on the corresponding training portion.
These assignments keep student-training examples, language-quality items, target-recovery items, and recorded-response endpoints distinct.
Appendix~\ref{app:data-targets} reports the exact inclusion and split definitions \evd{E005}\evd{E006}\evd{E008}\evd{E016}\evd{E025}\evd{E030}.

The training targets and controls differ in both information source and comparison role.
In the recorded-response route, the target contains observed Tuckute \ac{fmri} responses, while the permuted control preserves the response values and their marginal distribution but changes their pairing with the sentences.
In the synthetic-response route, the exact TRIBE target is produced by a fixed text-to-\ac{fmri} generator and contains 20,484 cortical-surface coordinates \autocite{dascoli2026}.
Its frozen row twin applies one fixed Sattolo permutation to complete target rows, eliminating unchanged stimulus--target pairs while preserving the target values and covariance structure.
The saved block permutation provides only a sensitivity analysis because its construction leaves fixed points and can duplicate or omit rows.
The text-derived auxiliary target instead mean-pools representations from a frozen \ac{lm} and applies a seeded Gaussian projection to the same 20,484 dimensions.
By construction, it matches the TRIBE target in corpus, dimension, model lineage, loss form, and training schedule.
Covariance geometry, nuisance predictability, learnability, optimization pressure, and retained-student movement require separate empirical checks, while initial-loss and gradient-scale comparability cannot be reconstructed from the retained artifacts.
Appendix~\ref{app:data-targets} specifies the exact generation, permutation, standardization, and dimensionality-reduction procedures \evd{E008}\evd{E016}\evd{E026}\evd{E030}.

The saved-student transfer analysis evaluates biological generalization after synthetic-response training.
Six saved students per training arm remain frozen throughout this analysis.
For each of the nine valid Tuckute participants, a fresh ridge readout is fitted within each of the same five contiguous folds, using 800 sentences for fitting and 200 for evaluation.
Layer~7 is the primary representation layer, and layer~6 provides a prespecified sensitivity analysis.
Student training and target-recovery evaluation use WikiText-derived examples, while the recorded Tuckute responses enter only the fold-specific evaluation readout.
This separation keeps the biological endpoint independent of student optimization and excludes reuse of the temporary training head.
Appendix~\ref{app:evidence-provenance} identifies the retained student artifacts, and Appendix~\ref{app:data-targets} specifies participant inclusion and response preprocessing \evd{E025}.

\subsection{Student training and retained components}
\label{sec:training}

The teacher is frozen and evaluated outside gradient computation.
On nonpadding token positions, logits are softened as $p_T^\tau=\operatorname{softmax}(z_T/\tau)$ and $q_S^\tau=\operatorname{softmax}(z_S/\tau)$.
The optimized objective is \(\mathcal L=\lambda_{\mathrm{text}}\mathcal L_{\mathrm{text}}+\lambda_{\mathrm{target}}\mathcal L_{\mathrm{target}}\), with \(\lambda_{\mathrm{text}}=1\).
For distillation arms, \(\mathcal L_{\mathrm{text}}\) is the attention-mask average of \(\tau^2D_{\mathrm{KL}}(p_T^\tau\Vert q_S^\tau)\); for squared-error target variants, \(\mathcal L_{\mathrm{target}}=\operatorname{MSE}(W\bar h_k,y)\), where \(\bar h_k\) is the attention-mask mean of the selected student state and \(W\) is a temporary linear head.
The \ac{kl} direction is therefore teacher to student; exact temperatures, weights, schedules, and early cross-entropy variants are reported in Appendix~\ref{app:training} \evd{E003}\evd{E004}\evd{E005}\evd{E016}.
Masking removes padded positions from both the token loss and the pooled state, and the $\tau^2$ factor retains the implemented temperature scaling.
The squared-error target loss is averaged over batch and target coordinates; it is a training signal, not the estimator used for the reported endpoint.

The text loss traverses the permanent \ac{lm} path.
The target loss branches at $h_k$ and updates the temporary head plus every trainable parameter that produces that state, including the tied embedding/output tensor through its input-embedding role in full-parameter TRIBE runs; it supplies no gradient through later blocks or the output-projection operation along this branch.
Those later blocks move only through the text loss, while the tied tensor also receives text-loss gradient through its output role.
Figure~\ref{fig:training-architecture} separates the two paths.

\begin{figure}[t]
  \centering
  \resizebox{\textwidth}{!}{%
  \begin{tikzpicture}[
    block/.style={draw,rounded corners,align=center,minimum height=8mm,text width=20mm,fill=blue!5},
    datum/.style={draw,rounded corners,align=center,minimum height=7mm,text width=15mm,fill=gray!8},
    loss/.style={draw,rounded corners,align=center,minimum height=8mm,text width=26mm,fill=purple!8},
    temp/.style={draw,dashed,rounded corners,align=center,minimum height=8mm,text width=21mm,fill=orange!10},
    forward/.style={-{Latex[length=2mm]},thick},
    grad/.style={-{Latex[length=2mm]},thick,dashed,red!70!black},
    panel/.style={draw=gray!60,rounded corners,inner sep=7pt}
  ]
    \node[datum] (text) {training text};
    \node[block,right=3mm of text] (embed) {retained input embedding};
    \node[block,right=3mm of embed] (early) {retained blocks $1{:}k$\\trainable or frozen by regime};
    \node[datum,right=3mm of early,fill=blue!5] (state) {middle state $h_k$};
    \node[block,right=3mm of state] (late) {retained blocks $k{+}1{:}L$};
    \node[block,right=3mm of late] (output) {retained tied \ac{lm} output $q_{\mathrm{student}}$};
    \draw[forward] (text)--(embed);
    \draw[forward] (embed)--(early);
    \draw[forward] (early)--(state);
    \draw[forward] (state)--(late);
    \draw[forward] (late)--(output);

    \node[datum,above=11mm of late] (teacher) {frozen teacher $p_{\mathrm{teacher}}$};
    \node[loss,above=11mm of output] (kd) {$\mathcal{L}_{\mathrm{KD}}$\\$\tau^2D_{\mathrm{KL}}(p_T^\tau\Vert q_S^\tau)$};
    \draw[forward] (text.north) |- (teacher.west);
    \draw[forward] (teacher)--(kd);
    \draw[forward] (output)--(kd);
    \draw[grad] (kd.north) -- ++(0,4mm) -| (embed.north);

    \node[temp,below=11mm of state] (nhead) {temporary target head\\discarded after training};
    \node[loss,right=5mm of nhead] (nloss) {$\mathcal{L}_{\mathrm{target}}$\\recorded or predicted target};
    \draw[forward] (state)--(nhead);
    \draw[forward] (nhead)--(nloss);
    \draw[grad] (nloss.south) -- ++(0,-4mm) -| (embed.south);

    \node[font=\bfseries,anchor=west] (training-title) at ($(text.north west)+(0,7mm)$) {Training};
    \node[panel,fit=(text)(embed)(early)(state)(late)(output)(teacher)(kd)(nhead)(nloss)(training-title)] {};

    \node[datum,below=58mm of text] (etext) {held-out text};
    \node[block,right=5mm of etext,text width=30mm] (frozen) {frozen retained student\\temporary head absent};
    \node[datum,right=5mm of frozen,fill=blue!5] (estate) {retained state $h_k$};
    \node[block,right=5mm of estate,text width=25mm,fill=green!7] (ridge) {fresh ridge readout\\fit inside evaluation folds};
    \node[datum,right=5mm of ridge,text width=22mm] (score) {held-out target\\$R^2$ or unique $R^2$};
    \draw[forward] (etext)--(frozen);
    \draw[forward] (frozen)--(estate);
    \draw[forward] (estate)--(ridge);
    \draw[forward] (ridge)--(score);
    \node[font=\bfseries,anchor=west] (evaluation-title) at ($(etext.north west)+(0,7mm)$) {Evaluation};
    \node[panel,fit=(etext)(frozen)(estate)(ridge)(score)(evaluation-title)] {};
  \end{tikzpicture}}
  \caption{Training and evaluation use different heads. Solid arrows show forward computation; red dashed arrows show gradient reach. The target branch reaches the temporary head and every trainable producer of the selected state, including embeddings when they are trainable, but not later blocks through that branch. Frozen parameters receive no updates. Evaluation freezes the retained student, discards the temporary head, and fits a new ridge readout.}
  \label{fig:training-architecture}
\end{figure}

The primary recorded-target regime distils Qwen2.5-1.5B into Qwen2.5-0.5B while training on the 800 Tuckute tuning sentences in each outer fold \evd{E005}\evd{E008}.
It uses attention-mask mean pooling at layer~12, five-dimensional \ac{roi} targets, three training seeds, and the same five outer folds; language quality is measured on held-out WikiText rather than on the neural evaluation items.
The averaged and participant-specific targets share this machinery but answer different questions \evd{E005}\evd{E008}.

The trainable subset is part of each intervention.
Primary recorded-target runs train \ac{lora} adapters \autocite{hu2022} and the temporary head while the base model and tied embedding/output matrix remain frozen; the naturalistic full-fine-tuning variant instead trains transformer weights and the temporary head while retaining that tied-matrix freeze \evd{E017}.
TRIBE trains all GPT-2 student parameters, including the tied embedding/output weights, plus the temporary head.
These are distinct parameterizations, so absolute scores across recorded and predicted families are not intervention contrasts \evd{E004}\evd{E008}\evd{E016}\evd{E017}.
Here ``retained'' means that a component survives deployment; it does not imply that the component was trainable in every regime.

Same-initialization, same-data, and same-budget statements apply only within a declared paired intervention family.
There, aligned, permuted, and ordinary-\ac{kd} arms share the seed, initialized student, text, optimizer schedule, and step budget while the intended supervision changes.
The cold-versus-warm headroom screen is an explicit exception rather than a paired intervention: its cold arm was randomly initialized and ran two epochs, whereas its pretrained warm arms ran one, so they differ in both initialization and optimizer-step count \evd{E003}.
Full schedules and secondary parameter settings are confined to Appendix~\ref{app:training}.

After training, the temporary head is discarded and the retained student is frozen.
Every held-out target or \ac{fmri} endpoint uses a newly fitted ridge readout \evd{E016}.
The participant-transfer and exact-substrate diagnostics load saved students or fixed targets without retraining \evd{E025}\evd{E030}.
For recorded-target sentence interventions, student training holds out one 200-sentence outer block; the fresh response readout is then fitted and evaluated wholly within that block using five contiguous inner subfolds.
For the baseline sentence and naturalistic-story assays, \texttt{RidgeCV} selects one penalty shared across all response coordinates from the relevant readout-training complement.
In every case, regularization and transforms are nested inside the fitting set, so the temporary training head and held-out responses cannot enter a reported score.

\FloatBarrier
\subsection{Measuring brain predictivity and target recovery}
\label{sec:measurement}

Ordinary held-out $R^2$ compares prediction error with variation around each coordinate's test-fold mean.
For target coordinate $j$ and test indices $\mathcal{T}$,
\[
R_j^2=1-\frac{\sum_{i\in\mathcal{T}}(y_{ij}-\widehat{y}_{ij})^2}
{\sum_{i\in\mathcal{T}}(y_{ij}-\bar{y}_{\mathcal{T}j})^2}.
\]
Negative scores are retained: they indicate predictions worse than the test-mean baseline, not missing values.

Controlled brain predictivity is the semipartial increment from contextual features beyond a frozen nuisance design of low-level and static language variables.
Under identical fold preprocessing and ridge selection, the reported quantity is
\[
R^2_{\mathrm{unique}}=R^2_{\mathrm{full}}-R^2_{\mathrm{nuisance}}.
\]
This predictive difference is not a causal allocation of biological variance \evd{E002}\evd{E006}.

The Tuckute and LeBel brain assays compute coordinatewise $R_j^2$ and then average the declared \acp{roi} or voxels and folds; the continuity target-recovery score uses the same equal-coordinate convention \evd{E016}.
The exact-substrate 20,484-coordinate target-recovery assay instead pools the held-out sum of squared errors (\(\mathrm{SSE}\)) and held-out total sum of squares (\(\mathrm{SST}\)) over all coordinates before computing \(1-\mathrm{SSE}/\mathrm{SST}\) within each fold, which weights coordinates by held-out variance; its target-to-response and composed student-to-target-to-response assays remain coordinatewise \evd{E030}.
Coordinates, folds, and seeds are never promoted to biological replicates \evd{E016}\evd{E030}.

Every encoding endpoint fits a fresh ridge estimator on frozen features.
Every learned operation is fitted only on the relevant readout-training complement before application to its held-out block.
Sentence analyses preserve contiguous item blocks; naturalistic analyses hold out whole stories and learn only \ac{pca}, assembled-design scaling, response centering, ridge selection, and the readout after the fixed per-story pipeline described above.
The baseline screen reports best layers descriptively within a fixed grid; later recorded-response assays prospectively reuse those choices, while the exact-substrate retention and student-to-target-to-response assays lock layer~6 as primary and layer~7 as non-rescuing sensitivity.
Reliability filters are fixed independently of evaluation outcomes.
The core sentence nuisance design contains token length, normalized item position, and a fixed static embedding; the exact-substrate primary design also includes imageability \evd{E030}.
The naturalistic design contains word and phoneme rate, duration, length, log frequency, and a static embedding; every compared arm reuses the same fixed per-story temporal arrays.
Contextual and high-dimensional static blocks enter at a declared capacity-matched rank.
These nuisance arrays are held identical within a contrast, so a training arm cannot gain merely from a wider readout or a different low-level timing model.
Appendices~\ref{app:data-targets} and~\ref{app:inference} specify nuisance blocks, ranks, temporal transforms, and secondary estimators \evd{E001}\evd{E002}\evd{E006}\evd{E025}\evd{E030}.

\subsection{Estimands, controls, quality criteria, and inference}
\label{sec:estimands}

Quality-aware headroom is measured before the intervention claims.
At the declared layers, the \texttt{gpt2-medium} teacher at layer~14 is compared with random and pretrained \texttt{gpt2} students at layer~7, before and after ordinary \ac{kd}; alignment is interpreted jointly with held-out perplexity rather than as an objective-only effect \evd{E003}.
The cross-tokenizer quality analysis uses tokenizer-comparable bits per byte \evd{E015}.

The recorded-target pairing effect is the fresh held-out response-predictivity difference between the recorded-target arm and its matched permuted-target arm.
Both arms use \(\lambda_{\mathrm{target}}=10\), the same seed, checkpoint schedule, optimizer-step budget, and fixed WikiText quality probe, and report held-out perplexity.
Comparisons with ordinary \ac{kd} require checkpoint or curve matching on perplexity; without that match they remain descriptive.
The averaged-response analysis remains a fixed-stimulus, fold-level estimand, whereas participant analyses aggregate folds and seeds within person before population inference \evd{E004}\evd{E005}\evd{E008}\evd{E017}.

The predicted branch follows claim order.
First, exact-target measurability asks whether the fixed TRIBE target predicts recorded responses beyond nuisance and its frozen row twin.
Second, target uptake and retained-student movement compare fresh held-out target recovery and post-training model change with seed-matched \ac{kd} and target-permutation arms.
Third, synthetic-response attribution asks whether TRIBE and projected text features are adequately matched in geometry, learnability, and induced movement and whether the brain-response-target arm outperforms that control on the claimed outcome.
Matching makes attribution to brain-response content testable; adequacy alone does not identify that attribution.
Only then does biological transfer compare saved target-trained students on participant-specific recorded responses; none of the earlier endpoints substitutes for this one \evd{E016}\evd{E025}\evd{E026}\evd{E030}.

Quality criteria are endpoint-specific.
The scaled predicted-target family measures perplexity on its exact 1,999-sentence held-out corpus and requires each target arm to remain within 5\% relative perplexity of its seed-matched \ac{kd} arm \evd{E016}.
The saved-student participant-transfer analysis rescored the same sentences as tokenizer-comparable bits per byte and requires $|\Delta\mathrm{bpb}|\leq0.005$ within seed \evd{E025}.
Perplexity is used only within a shared tokenizer and scoring protocol; cross-family quality comparisons use bits per byte.
If a required quality criterion or manipulation check fails, downstream scores are descriptive rather than claim-bearing \evd{E005}\evd{E016}\evd{E025}.

Participants are the sole population units for claims about people; story blocks, stimulus folds, voxels, coordinates, and seeds remain nested or technical measurements.
Uncertainty is therefore aligned to the stated unit, multiplicity is controlled across declared primary contrasts, and study-specific thresholds govern continuation without becoming universal biological cutoffs.
An external utility claim requires its own quality-matched task endpoint; a branch stopped at a precondition supplies no utility evidence.
The Results therefore follow the canonical evidence order: controlled brain predictivity, quality-aware headroom, recorded-target pairing, exact-target measurability, target uptake and retained-student movement, brain-response-specific attribution, biological transfer, and external utility.
